% -----------------------------------------------------------------------------
% This file is automatically generated.
% Do not edit manually.
% Source: notebooks/support/hypothesis-testing/support.Rmd
% -----------------------------------------------------------------------------


\noindent We start by setting up the MUS object \ttblue{ar} (for Accounts Receivable) in \ttblue{R} with the \ttblue{FSaudit} package and first verify that the number of sampling units \ttblue{popn} and the total book value \ttblue{popBv} of the sampling frame match those in the population. Notice that these statistics are calculated as soon as the object is loaded with the detail amounts.

\par\addvspace{\topsep}
\begingroup
\fvset{listparameters={%
  \setlength{\topsep}{0pt}%
  \setlength{\partopsep}{0pt}%
  \setlength{\parsep}{0pt}%
  \setlength{\itemsep}{0pt}%
}}
\begin{Verbatim}[breaklines=true,formatcom=\color{ada_blue}]
ar <- mus_obj(
  bv = accounts_receivable$amount,
  id = accounts_receivable$invoice
)
ar$popn
\end{Verbatim}
\begin{Verbatim}[breaklines=true,formatcom=\color{ada_light_blue}]
[1] 10000
\end{Verbatim}
\begin{Verbatim}[breaklines=true,formatcom=\color{ada_blue}]
ar$popBv
\end{Verbatim}
\begin{Verbatim}[breaklines=true,formatcom=\color{ada_light_blue}]
[1] 13500000
\end{Verbatim}
\endgroup
\par\addvspace{\topsep}
\noindent Sample size calculation is invoked with the relevant values from the \emph{Case: Accounts receivable circularization}.
